% Options for packages loaded elsewhere
\PassOptionsToPackage{unicode}{hyperref}
\PassOptionsToPackage{hyphens}{url}
\documentclass[
]{article}
\usepackage{xcolor}
\usepackage{amsmath,amssymb}
\setcounter{secnumdepth}{-\maxdimen} % remove section numbering
\usepackage{iftex}
\ifPDFTeX
  \usepackage[T1]{fontenc}
  \usepackage[utf8]{inputenc}
  \usepackage{textcomp} % provide euro and other symbols
\else % if luatex or xetex
  \usepackage{unicode-math} % this also loads fontspec
  \defaultfontfeatures{Scale=MatchLowercase}
  \defaultfontfeatures[\rmfamily]{Ligatures=TeX,Scale=1}
\fi
\usepackage{lmodern}
\ifPDFTeX\else
  % xetex/luatex font selection
\fi
% Use upquote if available, for straight quotes in verbatim environments
\IfFileExists{upquote.sty}{\usepackage{upquote}}{}
\IfFileExists{microtype.sty}{% use microtype if available
  \usepackage[]{microtype}
  \UseMicrotypeSet[protrusion]{basicmath} % disable protrusion for tt fonts
}{}
\makeatletter
\@ifundefined{KOMAClassName}{% if non-KOMA class
  \IfFileExists{parskip.sty}{%
    \usepackage{parskip}
  }{% else
    \setlength{\parindent}{0pt}
    \setlength{\parskip}{6pt plus 2pt minus 1pt}}
}{% if KOMA class
  \KOMAoptions{parskip=half}}
\makeatother
\usepackage{longtable,booktabs,array}
\newcounter{none} % for unnumbered tables
\usepackage{calc} % for calculating minipage widths
% Correct order of tables after \paragraph or \subparagraph
\usepackage{etoolbox}
\makeatletter
\patchcmd\longtable{\par}{\if@noskipsec\mbox{}\fi\par}{}{}
\makeatother
% Allow footnotes in longtable head/foot
\IfFileExists{footnotehyper.sty}{\usepackage{footnotehyper}}{\usepackage{footnote}}
\makesavenoteenv{longtable}
\usepackage{graphicx}
\makeatletter
\newsavebox\pandoc@box
\newcommand*\pandocbounded[1]{% scales image to fit in text height/width
  \sbox\pandoc@box{#1}%
  \Gscale@div\@tempa{\textheight}{\dimexpr\ht\pandoc@box+\dp\pandoc@box\relax}%
  \Gscale@div\@tempb{\linewidth}{\wd\pandoc@box}%
  \ifdim\@tempb\p@<\@tempa\p@\let\@tempa\@tempb\fi% select the smaller of both
  \ifdim\@tempa\p@<\p@\scalebox{\@tempa}{\usebox\pandoc@box}%
  \else\usebox{\pandoc@box}%
  \fi%
}
% Set default figure placement to htbp
\def\fps@figure{htbp}
\makeatother
\setlength{\emergencystretch}{3em} % prevent overfull lines
\providecommand{\tightlist}{%
  \setlength{\itemsep}{0pt}\setlength{\parskip}{0pt}}
% ---------------------------------------------------------------------------
% Unicode -> LaTeX mappings so these documents compile and render correctly
% under xelatex (and pdflatex) without depending on a particular system font.
% Every non-ASCII glyph used in the course text is mapped explicitly; this is
% needed because the documents use mathematical symbols inside running text
% and inside inline-code spans, where unicode-math would otherwise treat them
% as math-only and abort.
% ---------------------------------------------------------------------------
\usepackage{amssymb}
\usepackage{newunicodechar}

% --- Greek (lower case) ---
\newunicodechar{α}{\ensuremath{\alpha}}
\newunicodechar{β}{\ensuremath{\beta}}
\newunicodechar{γ}{\ensuremath{\gamma}}
\newunicodechar{δ}{\ensuremath{\delta}}
\newunicodechar{θ}{\ensuremath{\theta}}
\newunicodechar{κ}{\ensuremath{\kappa}}
\newunicodechar{λ}{\ensuremath{\lambda}}
\newunicodechar{μ}{\ensuremath{\mu}}
\newunicodechar{ν}{\ensuremath{\nu}}
\newunicodechar{π}{\ensuremath{\pi}}
\newunicodechar{ρ}{\ensuremath{\rho}}
\newunicodechar{ς}{\ensuremath{\varsigma}}
\newunicodechar{σ}{\ensuremath{\sigma}}
\newunicodechar{τ}{\ensuremath{\tau}}
\newunicodechar{φ}{\ensuremath{\varphi}}
\newunicodechar{χ}{\ensuremath{\chi}}
% --- Greek (capitals) ---
\newunicodechar{Γ}{\ensuremath{\Gamma}}
\newunicodechar{Δ}{\ensuremath{\Delta}}
\newunicodechar{Λ}{\ensuremath{\Lambda}}
\newunicodechar{Σ}{\ensuremath{\Sigma}}

% --- Math symbols ---
\newunicodechar{ℝ}{\ensuremath{\mathbb{R}}}
\newunicodechar{∂}{\ensuremath{\partial}}
\newunicodechar{∈}{\ensuremath{\in}}
\newunicodechar{∫}{\ensuremath{\int}}
\newunicodechar{√}{\ensuremath{\surd}}
\newunicodechar{≈}{\ensuremath{\approx}}
\newunicodechar{≤}{\ensuremath{\leq}}
\newunicodechar{≥}{\ensuremath{\geq}}
\newunicodechar{±}{\ensuremath{\pm}}
\newunicodechar{×}{\ensuremath{\times}}
\newunicodechar{−}{\ensuremath{-}}
\newunicodechar{·}{\ensuremath{\cdot}}
\newunicodechar{½}{\textonehalf}
\newunicodechar{→}{\ensuremath{\rightarrow}}
\newunicodechar{⟶}{\ensuremath{\longrightarrow}}

% --- Superscripts / modifier letters ---
\newunicodechar{²}{\ensuremath{^{2}}}
\newunicodechar{³}{\ensuremath{^{3}}}
\newunicodechar{⁺}{\ensuremath{^{+}}}
\newunicodechar{⁻}{\ensuremath{^{-}}}
\newunicodechar{ᵀ}{\ensuremath{^{\mathsf{T}}}}
\newunicodechar{ᵏ}{\ensuremath{^{k}}}

% --- Punctuation / misc text symbols ---
\newunicodechar{…}{\textellipsis}
\newunicodechar{§}{\S}
\newunicodechar{⚑}{\textbf{[!]}}

% --- Accented Latin (author names) ---
\newunicodechar{Á}{\'A}
\newunicodechar{Ø}{\O}
\newunicodechar{ä}{\"a}
\newunicodechar{é}{\'e}
\newunicodechar{ó}{\'o}
\newunicodechar{ô}{\^o}
\usepackage{bookmark}
\IfFileExists{xurl.sty}{\usepackage{xurl}}{} % add URL line breaks if available
\urlstyle{same}
\hypersetup{
  pdftitle={Tutor Handbook},
  pdfauthor={North Carolina State University · Financial Mathematics},
  hidelinks,
  pdfcreator={LaTeX via pandoc}}

\title{Tutor Handbook}
\usepackage{etoolbox}
\makeatletter
\providecommand{\subtitle}[1]{% add subtitle to \maketitle
  \apptocmd{\@title}{\par {\large #1 \par}}{}{}
}
\makeatother
\subtitle{Optimal Market Making for Cryptocurrency Options --- The
Avellaneda--Stoikov Framework and its Extension to Options}
\author{North Carolina State University · Financial Mathematics}
\date{Summer 2026}

\begin{document}
\maketitle

\includegraphics[width=2.6in,height=\textheight,keepaspectratio]{assets/ncsu_logo.png}

\section{Tutor Handbook}\label{tutor-handbook}

\textbf{Optimal Market Making for Cryptocurrency Options}

\emph{The Avellaneda--Stoikov Framework and its Extension to Options}

A 10-Week Summer Research Program --- North Carolina State University

Prepared for teaching assistants, recitation leaders, and graders.

\emph{Program Mentor: Dendi Suhubdy (Bitwyre). Guest Practitioner: Fenni
Kang (Coincall). Edition: 2026.}

\begin{center}\rule{0.5\linewidth}{0.5pt}\end{center}

\subsection{How to Use This Handbook}\label{how-to-use-this-handbook}

This handbook is the operational companion to the student-facing
syllabus. Read the syllabus first. The handbook has four parts: Part I
(orientation --- your role, the weekly rhythm, policies), Part II (the
week-by-week teaching guide), Part III (milestone evaluation rubrics),
and Part IV (administrative reference). The detailed per-problem
solution sketches and grading rubrics now live in the separate
\textbf{Solutions} booklet; this handbook tells you how to \emph{run}
each week and how to grade consistently, and points you to the Solutions
booklet for the worked answers.

\textbf{Conventions.} Boldface marks the single most important point in
a section. A ⚑ marks situations to escalate to the mentor. Time
allocations are recommendations; adapt to the cohort.

\textbf{A note on the two languages.} Problem sets are written in
\textbf{Python} with \textbf{PyTorch}; differentiation is done by
autograd, and students are \emph{not} asked to hand-derive AD. The
fast-computation week (Week 7) is written in \textbf{C++} and exposed to
Python via a \textbf{\texttt{pybind11}} binding. When you grade or help,
hold students to ``PyTorch for research and Greeks, C++ for the hot
path, \texttt{pybind11} between them.'' If a student hand-rolls adjoint
AD in C++, that is \emph{not} required and usually a sign they have
over-engineered --- steer them back to autograd as the reference.

\begin{center}\rule{0.5\linewidth}{0.5pt}\end{center}

\subsection{Part I. Tutor Orientation}\label{part-i.-tutor-orientation}

\subsubsection{1.1 Your Role}\label{your-role}

Tutors deliver the day-to-day program: leading the weekly ninety-minute
recitation, holding three weekly office hours, grading weekly problem
sets within seven days, giving first-pass evaluation of milestone
deliverables (final evaluation rests with the mentor), and flagging
at-risk students to the mentor within twenty-four hours. Tutors do
\textbf{not} deliver the main lectures or unilaterally change the
syllabus, problem sets, or rubrics --- raise proposed changes with the
mentor.

\subsubsection{1.2 Relationship with the
Mentor}\label{relationship-with-the-mentor}

The program is mentored by Dendi Suhubdy, who sets direction, delivers
the main lectures, and makes final decisions on grading disputes and
milestone evaluations. Standing channels: a weekly thirty-minute
tutor--mentor sync; a shared Slack workspace (mentor responds within 24
business hours); and a direct phone line for crises (⚑ student in acute
distress, suspected integrity violation, data-security incident) --- use
sparingly.

Note that two lectures are delivered by a \textbf{guest practitioner},
Fenni Kang (Chief Strategy Officer, Coincall): the taker-strategies
session in Week 9 and the structured-products session in Week 10.
Coordinate logistics for these with the mentor; you support them in
recitation as usual.

\subsubsection{1.3 Weekly Rhythm}\label{weekly-rhythm}

For a cohort of six to eight students, a typical week is: Sunday evening
prep (read this handbook's Part II section, \textasciitilde2--3 h);
attend Lecture 1 (Mon) and Lecture 2 (Wed) to track confusions; office
hours (Tue, Fri, Sat mornings); lead recitation (Thu); grade the
previous week's problem set over the weekend (\textasciitilde4--6 h).
Total ≈ 15--20 h/week, peaking in milestone weeks (5, 7, 9), when you
also join code review.

\subsubsection{1.4 Office Hours Guidance}\label{office-hours-guidance}

\begin{itemize}
\tightlist
\item
  \textbf{Resist solving problems for students.} Ask what they tried,
  ask what the next step is, give the smallest unblocking hint.
\item
  Distinguish conceptual confusion (return to lecture) from
  implementation confusion (read the code together).
\item
  The same student with the same confusion every week is a signal ---
  flag it to the mentor.
\item
  Office hours are not milestone code review. Do not debug milestone
  code in office hours; that happens in scheduled code-review sessions
  during milestone weeks, so your grading judgment stays independent.
\end{itemize}

\subsubsection{1.5 Grading Standards}\label{grading-standards}

Consistency matters more than the calibration of any single grade. With
two tutors, grade the first problem set's first three problems together
to calibrate, and spot-check each other thereafter.

Every problem set is a programming assignment. Specific standards:

\begin{itemize}
\tightlist
\item
  \textbf{Code is graded on correctness, reproducibility, and code
  quality, in that order.} Code that runs and produces the right answer
  with poor style is graded above code that does not run.
\item
  \textbf{Numerical-verification problems} (confirm an analytical result
  by computation) are graded on whether the verification is set up
  correctly and the reported error is the right order of magnitude ---
  not on a written proof.
\item
  \textbf{Short written interpretations} (typically printed by the code
  as a summary of computed results) are graded on substantive engagement
  with the actual numbers; restating the question earns zero.
\item
  For Week 7, additionally grade the \textbf{C++/\texttt{pybind11}
  boundary}: does the module build with CMake, import cleanly, and agree
  with the PyTorch reference Greeks?
\end{itemize}

Return graded work within seven days. Provide written feedback on at
least one problem per set, even at full marks. The worked solutions are
in the Solutions booklet; develop your own complete solution before
grading.

\subsubsection{1.6 Tutor Calibration
Session}\label{tutor-calibration-session}

Before the program, the mentor holds a half-day calibration session:
review the syllabus, walk through one full week of these notes, grade a
sample problem set as a group, and review the milestone rubrics.
Attendance is mandatory.

\begin{center}\rule{0.5\linewidth}{0.5pt}\end{center}

\subsection{Part II. Week-by-Week Teaching
Notes}\label{part-ii.-week-by-week-teaching-notes}

Each week section gives the spine of the week, lecture-emphasis notes,
common student confusions, and a recitation plan. Read the upcoming week
in full during the Sunday prep block. Grading rubrics and worked answers
are in the Solutions booklet.

\subsubsection{Week 1 --- Introduction
(Delivered)}\label{week-1-introduction-delivered}

This week was delivered last week. It oriented the cohort, set up the
Python + PyTorch and C++17/CMake/\texttt{pybind11} toolchains,
distributed Coincall data credentials, and collected the data-use
agreements. There is no graded problem set. If you are onboarding
mid-stream, confirm every student can (a) run a PyTorch autograd ``hello
world'' and (b) build and import the \texttt{pybind11} starter module
--- these are prerequisites for Weeks 5 and 7 respectively.

\subsubsection{Week 2 --- Foundations of Market
Microstructure}\label{week-2-foundations-of-market-microstructure}

\textbf{THE SPINE.} Students must leave able to think of a market maker
as a continuously quoting participant justified by the spread, and able
to read a limit order book without looking up vocabulary. Everything
builds on this.

\textbf{Lecture emphasis.} Lecture 1 is vocabulary and institutional
fluency; the key under-stated point is that price-time priority creates
a queue and queue position is a real economic asset (this matters in
Week 5 fill calibration). Lecture 2 covers Roll, Glosten--Milgrom, and
Kyle quickly --- the Glosten--Milgrom decision tree drawn on the board
is the most important artifact; have students copy it. Kyle is for
completeness.

\textbf{Common confusions.} Why quote both sides rather than react? (A
maker commits in advance; the spread compensates for that.) Aren't all
orders limit orders? (Mathematically yes, institutionally no ---
different fees and priority.) The Roll estimator differs from the
observed spread --- is that a bug? (No --- it estimates the unobserved
component; the gap is the finding.)

\textbf{Recitation.} Goal: every student leaves with a working Python
class that reconstructs an L2 book from an event stream. Materials: a
one-hour Coincall BTC-PERP event file (JSON), a tutor-only reference
solution, a projector. Timing: greet/recap (0--10); live demo of loading
the data and one minute of manual reconstruction (10--20); pair
programming on the scaffold (20--60) --- the most common bug is failing
to clean up empty price levels after a cancel; reconvene and debug one
pair's code together (60--75); wrap up lecture-flagged questions
(75--90). Save the working code; it is reused all program.

\subsubsection{Week 3 --- The Avellaneda--Stoikov
Model}\label{week-3-the-avellanedastoikov-model}

\textbf{THE SPINE.} Students must leave with the full stochastic-control
pipeline as a worked example: dynamics → HJB → ansatz → ODE reduction →
closed-form solution. Weeks 4--5 are variations on this template.

\textbf{Lecture emphasis.} Lecture 1's two key modelling choices:
arithmetic (not geometric) Brownian motion for the mid, and exponential
(CARA) utility, both for tractability --- make sure students see these
as choices. Lecture 2's centerpiece is the separation ansatz V =
−exp(−γ(x+qS))·θ(t,q); it recurs in every later control problem, so
attendance matters most here.

\textbf{Common confusions.} Why CARA? (Tractability --- it decouples the
value function from wealth.) The ansatz seems to come from nowhere (it
is principled --- multiplicative structure under CARA + arithmetic
dynamics; accept on faith, it recurs). Why is the half-spread roughly
symmetric but the reservation price asymmetric? (Different quantities
--- the reservation price is the indifference point; the half-spread is
the band width.) Huge simulated P\&L? (The model assumes you are alone
with exogenous flow; Weeks 5 and 9 add realism.)

\textbf{Recitation.} Group whiteboard derivation (0--30) --- the single
most important hour for many students; take your time. Then implement
the closed-form quotes in PyTorch (30--60) and run a γ sensitivity sweep
(60--80), aggregating inventory variance and Sharpe on the projector;
discuss where you would set γ for a real desk (80--90).

\subsubsection{Week 4 --- Stochastic Optimal Control in Continuous
Time}\label{week-4-stochastic-optimal-control-in-continuous-time}

\textbf{THE SPINE.} Students leave with the abstract template ---
dynamic programming principle, HJB, verification theorem --- and the
ability to recognize when a problem fits it. The Merton problem
reinforces, not introduces, the Week 3 pattern.

\textbf{Lecture emphasis.} Lecture 1 derives HJB in generality with
Avellaneda--Stoikov as a special case; the infinitesimal generator often
needs re-explaining for weaker stochastic-calculus backgrounds. Lecture
2 states the verification theorem (the missing Week-3 piece) and
contrasts CARA with CRRA. Do not let students get stuck on viscosity
solutions --- not central here.

\textbf{Common confusions.} Why is HJB necessary? (It encodes the DPP,
which is necessary by definition; the verification theorem gives
sufficiency.) Why does wealth appear under CRRA but not CARA? (Constant
marginal utility of wealth under CARA.) Verification vs.~existence
theorem? (We always have a candidate, so we only need verification.)

\textbf{Recitation.} Derive Merton under CARA (10--45), then CRRA
(45--75), faster the second time; side-by-side comparison (75--90). If
the cohort is strong, compress CARA and spend time on Almgren--Chriss
instead.

\subsubsection{Week 5 --- The Avellaneda--Stoikov Engine ⟶ Milestone
1}\label{week-5-the-avellanedastoikov-engine-milestone-1}

\textbf{THE SPINE.} Two things: students understand how the closed form
becomes running code, and they submit Milestone 1, the linear-asset
Avellaneda--Stoikov engine in Python/PyTorch.

\textbf{Lecture emphasis.} Lecture 1 is architecture and the move from
continuous time to discrete events; emphasize calibrating
λ(δ)=A·exp(−κδ) by maximum likelihood rather than a naive curve fit.
Lecture 2 is calibration and validation, and formally announces the
milestone and rubric.

\textbf{Common confusions.} Calibrated intensities look implausible
(crypto fill data is noisy; fit on one week, then two, then four ---
estimates should stabilize). Backtest P\&L is negative (check: spread
too tight vs. calibrated intensities, inventory penalty too small, or a
regime shift in the window --- before suspecting the model). My PyTorch
engine is slow (vectorize over inventory; avoid Python loops over the
grid).

\textbf{Recitation.} ⚑ Milestone week --- block four extra hours of code
review. Whiteboard any remaining derivation gaps (0--25), then live code
review of one volunteer's engine (25--70) --- this reveals cohort-wide
bugs you address in the wrap-up --- then open Q\&A on the milestone
acceptance criteria (70--90).

\subsubsection{Week 6 --- Options Pricing, the Volatility Surface, and
Calibration}\label{week-6-options-pricing-the-volatility-surface-and-calibration}

\textbf{THE SPINE.} Two tracks: conceptually, implied-vol surfaces and
their no-arbitrage constraints; practically, calibrating SVI and SABR to
messy Coincall data. Students should leave knowing calibration is half
craft, half theory.

\textbf{Lecture emphasis.} Lecture 1 spends the most time on calendar
and butterfly arbitrage (the constraints naive parameterizations
violate); vanna/volga are introduced but not tested until Week 8.
Lecture 2 live-demos a calibration that first fails --- the diagnostic
process is more valuable than the eventual fit. Students fit using
PyTorch autograd for the Jacobians, so they need not derive gradients by
hand.

\textbf{Common confusions.} Why two parameterizations? (They encode
different priors and fail in different regimes.) Huge residuals on
far-OTM puts? (Wide spreads / low volume --- weight by volume or inverse
spread.) SVI params jump day to day? (Often the initial guess --- seed
from yesterday's fit.) SABR vs.~Monte Carlo mismatch at long maturities?
(Hagan is a small-time expansion; MC is the truth.)

\textbf{Recitation.} Every student produces a working SVI and SABR
calibration with diagnostics on a chosen snapshot. Keep a backup,
slightly messier snapshot ready in case the chosen one is too clean for
the ``crypto calibration is hard'' lesson to land.

\subsubsection{Week 7 --- Fast Computation (Prices, Greeks, Implied Vol,
Surfaces) ⟶ Milestone
2}\label{week-7-fast-computation-prices-greeks-implied-vol-surfaces-milestone-2}

\textbf{THE SPINE.} The week is about latency. Students must internalize
that a correct but slow pricer is not enough. The tools are
characteristic- function pricing, robust implied-vol inversion, and ---
the new emphasis --- \textbf{binding C++ to Python with
\texttt{pybind11}}. The end is a sub-5ms surface revaluation callable
from Python.

\textbf{Lecture emphasis.} Lecture 1: COS is the workhorse (Carr--Madan
for context); the truncation interval must be computed from the
cumulants, not guessed; robust implied-vol inversion needs a good
initial guess (Jäckel-style) plus a safeguarded Newton/Halley step.
Lecture 2 is the \textbf{binder lecture}: structure a \texttt{pybind11}
module, build it with CMake, pass NumPy/PyTorch tensors without copying
(buffer protocol / \texttt{py::array\_t}), release the GIL around the
hot loop, and validate C++ Greeks against PyTorch autograd.
\textbf{Emphasize: autograd is the oracle for Greeks; C++ is for speed.}
We deliberately avoid hand-written adjoint AD in C++ --- students found
it confusing and it is unnecessary when autograd provides the reference.

\textbf{Common confusions.} ``Why are CF methods faster than
Black--Scholes?'' (They're not for BS --- the point is Heston, where
there's no closed form.) ``My \texttt{pybind11} build can't find
Python/pybind11 headers.'' (CMake \texttt{find\_package(pybind11)}; use
the provided \texttt{CMakeLists.txt} scaffold.) ``My C++ Greeks disagree
with autograd.'' (Usually a units or finite-difference-bump issue in the
comparison harness --- or the GIL not released causing a benchmarking
artifact.) ``COS won't match within 1e-4.'' (Truncation interval from
cumulants; raise N to 2\^{}10+.)

\textbf{Recitation.} ⚑ Milestone week --- block four extra hours of code
review. Profile a naive Python pricer (10--40), move the kernel to C++
and bind it with \texttt{pybind11} (40--70), re-benchmark and confirm
the surface revalues under 5 ms with Greeks matching autograd (70--90).
Have the reference timings ready so students know what ``good'' looks
like. The latency target is a hard gate --- be ready for arguments and
refer to the rubric.

\subsubsection{Week 8 --- Options Market Making Theory and Dynamic
Hedging}\label{week-8-options-market-making-theory-and-dynamic-hedging}

\textbf{THE SPINE.} The conceptual centerpiece: options market making is
Avellaneda--Stoikov with a vector-valued inventory, and the quoting
strategy interacts with an explicit hedging policy. The most cognitively
demanding week.

\textbf{Lecture emphasis.} Lecture 1 replaces scalar inventory with a
vector of Greeks and the scalar penalty with a quadratic form, derives
the multi-Greek HJB, and verifies reduction to Avellaneda--Stoikov in
1-D. Discourage setting the penalty matrix to the identity (defer the
``how to choose it'' discussion to office hours). Lecture 2 (Taleb)
derives the gamma--theta--variance identity --- the conceptual key to
options trading --- and presents external vs.~integrated hedging and
perpetual futures as the hedge instrument.

\textbf{Common confusions.} Why is inventory a vector? (Aggregated
Greeks, not contract positions --- long one call and short another at a
different strike is zero contracts but nonzero gamma/vega.) How to
choose the penalty matrix? (Calibrate to the empirical Greek covariance
or a desired risk-equivalence; it is part of the model design.) Hedge
with perps or the underlying? (Perps --- lower cost and the options are
quoted against the perp index; account for funding.)

\textbf{Recitation.} Multi-Greek HJB derivation (0--35), then a PyTorch
hedged-book simulation under two policies (35--70) with P\&L
distributions aggregated on the projector; discuss which policy wins in
which regime and tie to the Milestone 3 design (70--90).

\subsubsection{Week 9 --- Backtesting, Risk Analysis, and Taker
Strategies ⟶ Milestone
3}\label{week-9-backtesting-risk-analysis-and-taker-strategies-milestone-3}

\textbf{THE SPINE.} Students assemble everything into a working backtest
on three months of Coincall data --- more engineering discipline than
new theory --- and gain a practitioner's view of taker flow from the
guest lecture.

\textbf{Lecture emphasis.} Lecture 1 covers event-driven architecture
and the pitfalls (look-ahead, unrealistic fills, expiry handling) and
the attribution sum-to-total check (a mismatch is always a bug).
\textbf{Lecture 2 is the guest session (Fenni Kang) on taker strategies}
--- when/why takers cross the spread, the signals behind aggressive
flow, and what taker behaviour implies for order-flow toxicity. Brief
students beforehand to come with one question linking taker flow to
their own fill simulator.

\textbf{Common confusions.} Suspiciously high P\&L? (Always a bug ---
look-ahead in the surface, 100\% fills, missing costs/funding.)
Fill-simulator calibration seems circular? (Use a holdout: calibrate on
one month, validate on a disjoint month.) Attribution doesn't sum to
total? (Timing --- streams computed from different state snapshots.)
Sharpe of 6 over three months? (Estimation error --- require a
block-bootstrap CI and the deflated Sharpe.)

\textbf{Recitation.} ⚑ Milestone week --- six to eight hours of code
review. Each student presents one regime (one week of the window) and
walks through the P\&L attribution; probe for the pitfalls. This is the
dress rehearsal for the milestone.

\subsubsection{Week 10 --- Advanced Avellaneda--Stoikov, Structured
Products, and Final
Presentations}\label{week-10-advanced-avellanedastoikov-structured-products-and-final-presentations}

\textbf{THE SPINE.} Students choose one advanced Avellaneda--Stoikov
extension, implement it in a week, and present the final report. Your
role shifts from instruction to mentorship: help each student scope
realistically --- a clean partial result beats an incomplete grand one.

\textbf{Lecture emphasis.} Lecture 1 surveys the extension menu
(multi-asset, adverse selection, signal-driven, discrete inventory,
robust); students choose by end of day. \textbf{Lecture 2 is the guest
session (Fenni Kang) on corporate usage of options and structured
products} --- collars, accumulators/decumulators, and coupon products
--- and how that flow lands on the market-maker's book. It closes with
research-communication guidance.

\textbf{Common confusions.} ``My extension isn't working --- should I
switch?'' (No time to restart --- write up the smallest version that
works and acknowledge limitations.) Report length? (≤ 30 pages, quality
over length.) LLM use? (Writing support only, not content generation ---
see Part IV.)

\textbf{Recitation.} Final dry-run presentations (15 minutes each),
cohort feedback, tutor focused on the questions an industry audience
would ask. Be tough but constructive --- expose problems while there is
still time to fix them. Submit feedback rubrics to the mentor.

\begin{center}\rule{0.5\linewidth}{0.5pt}\end{center}

\subsection{Part III. Milestone Evaluation
Rubrics}\label{part-iii.-milestone-evaluation-rubrics}

Each milestone is graded out of 100. The final milestone grade is the
average of two independent evaluations (yours and the mentor's);
discrepancies over 10 points trigger a third reviewer.

\subsubsection{3.1 Milestone 1 --- Linear-Asset Avellaneda--Stoikov
Engine (Week
5)}\label{milestone-1-linear-asset-avellanedastoikov-engine-week-5}

{\def\LTcaptype{none} % do not increment counter
\begin{longtable}[]{@{}
  >{\raggedright\arraybackslash}p{(\linewidth - 4\tabcolsep) * \real{0.3793}}
  >{\raggedleft\arraybackslash}p{(\linewidth - 4\tabcolsep) * \real{0.1724}}
  >{\raggedright\arraybackslash}p{(\linewidth - 4\tabcolsep) * \real{0.4483}}@{}}
\toprule\noalign{}
\begin{minipage}[b]{\linewidth}\raggedright
Criterion
\end{minipage} & \begin{minipage}[b]{\linewidth}\raggedleft
Pts
\end{minipage} & \begin{minipage}[b]{\linewidth}\raggedright
Description
\end{minipage} \\
\midrule\noalign{}
\endhead
\bottomrule\noalign{}
\endlastfoot
Correctness & 30 & PyTorch quotes match the closed-form asymptotic
formula within 1e-6 on a synthetic problem with known solution. \\
Intensity calibration & 20 & ML calibration with χ² goodness-of-fit;
student demonstrates awareness of the fit's limitations. \\
Reproducibility & 15 & Identical config → identical results under a
fixed seed. Non-reproducibility is an automatic −15. \\
Backtest match & 15 & P\&L within ±50\% of the mentor's reference on the
same window. \\
Vectorization & 10 & Quote computation vectorized over inventory in
PyTorch (no Python grid loop). \\
Code quality & 5 & Lints; type hints; \texttt{pytest} over the core
quote computation. \\
Report & 5 & ≤ 10 pp.; follows the Avellaneda--Stoikov paper structure;
honest limitations. \\
\end{longtable}
}

\subsubsection{3.2 Milestone 2 --- Fast Pricing and Greeks Library (Week
7)}\label{milestone-2-fast-pricing-and-greeks-library-week-7}

{\def\LTcaptype{none} % do not increment counter
\begin{longtable}[]{@{}
  >{\raggedright\arraybackslash}p{(\linewidth - 4\tabcolsep) * \real{0.3793}}
  >{\raggedleft\arraybackslash}p{(\linewidth - 4\tabcolsep) * \real{0.1724}}
  >{\raggedright\arraybackslash}p{(\linewidth - 4\tabcolsep) * \real{0.4483}}@{}}
\toprule\noalign{}
\begin{minipage}[b]{\linewidth}\raggedright
Criterion
\end{minipage} & \begin{minipage}[b]{\linewidth}\raggedleft
Pts
\end{minipage} & \begin{minipage}[b]{\linewidth}\raggedright
Description
\end{minipage} \\
\midrule\noalign{}
\endhead
\bottomrule\noalign{}
\endlastfoot
Pricing accuracy & 20 & Max relative error 1e-4 vs.~Coincall reference
across the BTC surface. \\
Greek accuracy & 15 & C++ Greeks within 1e-4 of the PyTorch-autograd
reference for all six Greeks (delta, gamma, vega, theta, vanna,
volga). \\
Latency (full surface) & 25 & 300--500 contracts revalued with full
Greek vector in under 5 ms, one CPU core. \textbf{Hard gate: failing
caps the grade at 70.} \\
\texttt{pybind11} integration & 10 & Module builds with CMake, imports
cleanly from Python, passes tensors without unnecessary copies. \\
Implied-vol robustness & 10 & Inversion converges across the surface
including deep wings and short expiries. \\
Tests & 10 & Cover all pricers, the implied-vol solver, and the surface
adapter, plus one edge case. \\
Technical note & 10 & Latency reported as a distribution (histogram),
with optimization choices explained. \\
\end{longtable}
}

\subsubsection{3.3 Milestone 3 --- Full Options Market-Making Backtest
(Week 9)}\label{milestone-3-full-options-market-making-backtest-week-9}

{\def\LTcaptype{none} % do not increment counter
\begin{longtable}[]{@{}
  >{\raggedright\arraybackslash}p{(\linewidth - 4\tabcolsep) * \real{0.3793}}
  >{\raggedleft\arraybackslash}p{(\linewidth - 4\tabcolsep) * \real{0.1724}}
  >{\raggedright\arraybackslash}p{(\linewidth - 4\tabcolsep) * \real{0.4483}}@{}}
\toprule\noalign{}
\begin{minipage}[b]{\linewidth}\raggedright
Criterion
\end{minipage} & \begin{minipage}[b]{\linewidth}\raggedleft
Pts
\end{minipage} & \begin{minipage}[b]{\linewidth}\raggedright
Description
\end{minipage} \\
\midrule\noalign{}
\endhead
\bottomrule\noalign{}
\endlastfoot
Reproducibility & 10 & Identical replay → identical results.
Non-reproducibility is an automatic −10. \\
P\&L attribution & 20 & Streams (spread, inventory, gamma, vega,
hedging) sum to total P\&L within 1 bp. \\
Risk metrics & 15 & Sharpe, max drawdown, Calmar with block-bootstrap
CIs. The CI requirement is non-negotiable. \\
Greek-exposure tracking & 10 & Aggregated Greek time series plotted;
exposures within configured limits. \\
Engine quality & 15 & Correctly integrates the Week-7 pricing library
(via \texttt{pybind11}) with the Week-5 engine. \\
Hedging integration & 10 & Dynamic hedge correctly integrated; funding
rate accounted for. \\
Forensic analysis & 10 & Drawdown forensic identifies specific
positions, the environment, and a counterfactual. \\
Report & 10 & ≤ 20 pp.; includes the forensic; honest failure-mode
discussion. \\
\end{longtable}
}

Common bugs to check: look-ahead in the volatility surface (−15 engine
quality), 100\% fill assumption (−10), funding ignored (−5 hedging),
Sharpe without CI (−8 risk).

\subsubsection{3.4 Final Report (25 pts) and Presentation (15
pts)}\label{final-report-25-pts-and-presentation-15-pts}

\textbf{Report:} problem motivation (3), methodology (6), results with
uncertainty quantification (6), discussion (4), an explicit limitations
section (3), writing quality (3). \textbf{Presentation:} structure (3),
technical depth and Q\&A handling (4 + 2), visual aids (3), delivery
(3).

\begin{center}\rule{0.5\linewidth}{0.5pt}\end{center}

\subsection{Part IV. Administrative
Reference}\label{part-iv.-administrative-reference}

\subsubsection{4.1 Academic Integrity}\label{academic-integrity}

Follows the NC State Code of Student Conduct. Collaboration on problem
sets is encouraged but submissions must be each student's own work
(discussing the approach is fine; copying code is not). Cite any
external source, including LLM-generated content. LLM assistance is
allowed for writing support (grammar, phrasing) but not for technical
content generation; disclose non-trivial use. ⚑ The mentor handles all
integrity cases --- document specifics and escalate within 24 hours; do
not confront the student.

\subsubsection{4.2 Late Work}\label{late-work}

Problem sets are due Sunday 11:59 pm Eastern; late until Tuesday 11:59
pm with a 20\% deduction; nothing accepted after. Milestones are due
Friday 11:59 pm Eastern of the milestone week; no late milestones.
Emergencies (documented medical/family) are handled by the mentor only.

\subsubsection{4.3 Accommodations}\label{accommodations}

Implement documented accommodations as filed with the NC State
Disability Resource Office; do not ask students to justify them. If
unsure how to implement one, ask the mentor.

\subsubsection{4.4 Crisis Escalation ⚑}\label{crisis-escalation}

Escalate immediately (use the phone line, not Slack): suspected
integrity violation; student in acute distress; data-security incident
(e.g., a student exposes Coincall data publicly); hostile interaction.
The mentor directs next steps.

\subsubsection{4.5 Environment and
Tooling}\label{environment-and-tooling}

Each tutor receives Slack and LMS access, Coincall S3 credentials (under
NDA, plus reference-implementation credentials), a desk in SAS Hall, and
a laptop with the standard program environment pre-installed:

\begin{itemize}
\tightlist
\item
  \textbf{Python 3.11}, \textbf{PyTorch}, NumPy, SciPy, Pandas, PyArrow,
  and the program's own internal Python libraries.
\item
  A \textbf{C++17 toolchain} (clang or g++), \textbf{CMake},
  \textbf{Eigen}, a C++ \textbf{FFT} library (FFTW or pocketfft), and
  \textbf{\texttt{pybind11}} --- used in Week 7 to build and import the
  fast-computation module from Python.
\item
  Catch2 / GoogleTest and \texttt{pytest} for the test suites.
\end{itemize}

First-week setup: complete the data-use agreement; verify Slack/LMS
access; build and import the Week-7 \texttt{pybind11} starter module end
to end; run the reference Milestone 1 implementation; attend the
calibration session.

\subsubsection{4.6 Compensation, Hours, and
Contacts}\label{compensation-hours-and-contacts}

Tutors are compensated as standard graduate TAs under the NC State
Mathematics pay scale; hours are tracked through the university system.
The contact list (mentor, coordinator, departmental administrator, IT,
emergency) is distributed at orientation --- keep it accessible.

\end{document}
